% 上下两部分是独立的教学时间范围；纵向位置不承载数值。
\begin{tikzpicture}[x=1cm,y=1cm,font=\figurefont\small]
  \node[anchor=west] at (0,3.7) {时间点预测左右边界};
  \draw[-{Stealth[length=2mm]},BookInk,line width=0.8pt]
    (1,2.45) -- (10,2.45) node[below left] {秒};
  \fill[BookTeal!14] (2,2.85) rectangle (8.4,3.2);
  \draw[BookTeal,line width=1pt] (2,2.7) -- (2,3.35);
  \draw[BookTeal,line width=1pt] (8.4,2.7) -- (8.4,3.35);
  \fill[BookInk] (4.4,3.02) circle (2pt);
  \draw[{Stealth[length=1.7mm]}-{Stealth[length=1.7mm]},BookInk,line width=0.8pt]
    (2,3.02) -- (4.4,3.02);
  \draw[{Stealth[length=1.7mm]}-{Stealth[length=1.7mm]},BookInk,line width=0.8pt]
    (4.4,3.02) -- (8.4,3.02);
  \node[above] at (3.2,3.1) {3秒};
  \node[above] at (6.4,3.1) {5秒};
  \foreach \x/\lab in {2/9,4.4/12,8.4/17}{
    \draw[BookInk,line width=0.7pt] (\x,2.4) -- (\x,2.5);
    \node[below] at (\x,2.4) {\lab};
  }
  \node[anchor=west] at (0,1.55) {逐帧标签与片段碎裂};
  \node[anchor=east] at (1.2,0.9) {真实};
  \node[anchor=east] at (1.2,0.2) {预测};
  \foreach \x/\gt/\pr in {1.5/拿/拿,2.8/倒/倒,4.1/倒/拿,5.4/倒/倒,6.7/放/放,8/放/放}{
    \draw[BookRule,fill=BookPaper,line width=0.5pt] (\x,0.62) rectangle ++(1.2,0.55);
    \node at (\x+0.6,0.9) {\gt};
    \draw[BookRule,fill=BookPaper,line width=0.5pt] (\x,-0.08) rectangle ++(1.2,0.55);
    \node at (\x+0.6,0.2) {\pr};
  }
  \draw[BookAmber,line width=1.1pt] (4.1,-0.08) rectangle (5.3,0.47);
\end{tikzpicture}
